
% Authoritative LaTeX manuscript; Markdown is a generated compatibility projection.

% <!-- part-title: Orientation and Reference -->

% <!-- chapter-title: What This Book Argues -->

\chapter{What This Book Argues}

\emph{This section is a compact statement of the book\textquotesingle s argument, intended for orientation and later reference. The chapters that follow develop and support these claims. This book makes six central claims. Together they explain how commodity intelligence changes the economics of engineering, why scale remains a problem, and how MAGE responds.}

MAGE is a framework for governing consequential agentic engineering work. This book develops that framework through software engineering, its first worked domain and the source of its evidence. The six claims below are therefore stated primarily in software-engineering terms; \magevalue{chapter:theory} abstracts their mechanisms, and \magevalue{chapter:profession} asks what transfers beyond software.

\begin{enumerate}
\def\labelenumi{\arabic{enumi}.}
\tightlist
\item
  \textbf{Commodity intelligence changes the economics of software engineering.} Implementation capacity is becoming abundant relative to engineering judgment. As implementation gets cheaper, engineering effort shifts toward what remains scarce: deciding what to build, representing the system clearly enough to reason about it, producing evidence, and making important requirements enforceable by the engineering environment.
\item
  \textbf{Scale creates a reasoning problem.} Large software systems already exceed the reasoning horizon of humans; agents inherit the same problem. Software engineering has always answered scale with abstraction. Commodity intelligence does not remove that need. It makes the representations that guide the work more important.
\item
  \textbf{Purposeful models can reduce the cost of engineering reasoning.} Software engineers have always reasoned about architecture and other system properties through abstraction. What changes is how much of that reasoning can economically remain explicit. Purposeful models capture the knowledge and intent needed for engineering decisions while leaving irrelevant choices open. As commodity intelligence lowers the cost of deriving, maintaining, and using such representations, more engineering knowledge can be carried forward rather than reconstructed from implementation. Models let humans and agents reason about larger properties while deliberately leaving implementation choices open where engineering has imposed no obligation.
\item
  \textbf{Explicit obligations can make autonomous implementation governable.} When agents perform more of the implementation, deciding what the system must do becomes increasingly separate from writing the code that does it. Important engineering decisions therefore cannot live only in instructions to an agent or in a person\textquotesingle s head. Alignment encodes selected obligations into checks and controls that can constrain work or determine what the environment will accept. This distinguishes choices an agent is free to make from decisions engineering has already made. Engineers can then give agents substantial freedom in how they build the system while retaining control over the properties that matter.
\item
  \textbf{Governance conversion turns recurring judgment into durable engineering structure.} When a failure exposes missing knowledge or an unenforced obligation, encode the lesson into a model, procedure, or mechanism that future work can inherit. Durable structure becomes engineering capital when later work keeps benefiting from it.
\item
  \textbf{Engineering work will reorganize around what remains scarce.} As implementation becomes cheaper, more engineering effort will move toward representation, evidence, governance, coordination, and judgment. Agents may perform increasing portions of that work as well. The enduring boundary will not be a particular task or level of abstraction, but responsibility for deciding what matters, what evidence is sufficient, which obligations should be enforced, and what tradeoffs remain acceptable.
\end{enumerate}

MAGE has two principles and a conversion loop. Modeling makes consequential knowledge and intent explicit, while Alignment makes selected obligations enforceable. Governance conversion turns recurring judgment into durable engineering structure that future work can inherit. \autoref{mage-method} connects these responses to the economic change, the reasoning problem, and the predicted reorganization of engineering work.

\begin{figure}[htbp]
\centering
\includegraphics[width=\linewidth]{book/latex/assets/mage-method-5b47654b89.pdf}
\caption{The six claims, connected. The feedback loop turns recurring judgment into structure that later work can inherit.}
\label{mage-method}
\end{figure}
